\scanpage{29}
\margin{28}

les pts de $\widetilde\Delta$ au dessus de $s$ i.e. les orientations de $X$ en $s$, correspondent aux \inserted{deux} vis de $\Delta\setminus\{s\}$ dans $X$. La spécialisation $D\mathbin{-}D'$ définit une de ces vis (celle balayée par la spécialisation), et on trouve un isom. combinatoire canonique entre les segments combinatoires correspondants : la vis de spécialisation sur $D$, et celle correspondant à \inserted{l'intérieur de} la vis balayée de $D'\setminus\{s\}$, \inserted{i.e. le \uline{\unsure{découpage}} de $D'$ \unsure{suivant} $s$.} (En termes de $D'$, le choix d'un \deleted{point} \inserted{sommet} $\widetilde s$ de $\widetilde{D'}$, \uline{d'où} un antipodique $\widetilde s^-$, partage $\widetilde{D'}$ en deux arcs, correspondant aux deux orientations de $\widetilde{D'}$ en les considérant comme \deleted{les} deux arcs \uline{orientés} d'origine $\widetilde s$, extrémité $\widetilde s^-$).

\deleted{L'antipodisme de $\widetilde{D'}$ qui \unsure{échange les deux} \inserted{échange $\widetilde s,\widetilde s^-$} \unsure{soit l'identité ou}}

\deleted{\unsure{Il y a un unique autre} \deleted{\unsure{isom}} combinatoire \inserted{$\ne\mathrm{id}$} \unsure{qui identifie} ces deux arcs, et qui \deleted{\unsure{les deux arcs}} --- par cet isom, on peut identifier les deux arcs combinatoires obtenus, d'où un arc combinatoire canonique.}

L'antipodisme définit un isom canonique entre ces deux arcs, qui permet \inserted{donc} d'identifier combinatoirement, \inserted{\unsure{soit} $\unsure{\mathbb{I}|\mathrm{I}}(D',s)$} \unsure{ce} segment combinatoire. D'où ce qu'on peut dire, c'est qu'on a un isom combinatoire canonique
\diagram{\begin{tikzcd}[ampersand replacement=\&]
\begin{gathered}\text{Vis de spécialisation}\\\text{de $D$ utilisée}\end{gathered}\arrow[r,"\sim"]\&\unsure{\mathbb{I}|\mathrm{I}}(D',s)
\end{tikzcd}}
